\chapter{بنى لويس والرنين ونموذج VSEPR}\label{ch:b1:lewis-resonance-vsepr}

جزيء الأوزون، \ce{O3}، ثلاث ذرات أكسجين في سلسلة منحنية. ارسمه بقواعد
الكتاب الأول تخرج إحدى الرابطتين ثنائية والأخرى أحادية: فينبغي أن تكون
إحداهما قصيرة والأخرى طويلة. غير أن مطيافية الموجات الدقيقة تجد رابطتي
\ce{O-O} في الأوزون متساويتين تمامًا، عند
\qty{127.8}{pm}، بين الرابطة الثنائية في \ce{O2} (\qty{120.8}{pm})
والرابطة الأحادية في بيروكسيد الهيدروجين (\qty{147.5}{pm}). فلا رسم واحد
صحيح؛ بل رسمان معًا. يجعل هذا الفصل نموذج لويس
دقيقًا --- \omterm{def:b1:lewis-resonance-vsepr:formal-charge}{الشحنات الشكلية}، والثمانيات الناقصة أو المتجاوَزة، والرنين
بين عدة بنى --- ثم يستعمله للتنبؤ بأشكال الجزيئات وبعزوم ثنائي القطب
الناتجة عنها.
% ledger: geo:O3.rOO, re:O2, geo:H2O2.rOO

\begin{recall}
وصف الكتاب الأول (الصف العاشر) \omterm{def:b1:lewis-resonance-vsepr:lewis}{الرابطة التساهمية} بأنها زوج إلكترونات
مشترك، ورسم \omterm{def:b1:lewis-resonance-vsepr:lewis}{بنى لويس} بأزواجها الرابطة والحرة، وتنبأ
بأربعة أشكال (خطي، منحنٍ، مثلثي، رباعي الأوجه). وقد صارت
\omterm{def:b1:periodicity:electronegativity}{الكهرسلبية} وفروقها كمّية في
\cref{ch:b1:periodicity}؛ وتُقرأ \omterm{def:b1:quantum-numbers:configuration}{إلكترونات التكافؤ} لذرة من
توزيعها (\cref{def:b1:quantum-numbers:configuration}).
\end{recall}

\section{بنى لويس والشحنات الشكلية}

\begin{definition}[الرابطة التساهمية، بنية لويس]\label{def:b1:lewis-resonance-vsepr:lewis}
\emph{الرابطة التساهمية}\index{رابطة تساهمية} زوج إلكترونات تتشاركه
ذرتان، هو \emph{الزوج الرابط}\index{زوج رابط}؛ ويكوّن زوجان أو ثلاثة
أزواج رابطة ثنائية أو ثلاثية. وزوج \omterm{def:b1:quantum-numbers:configuration}{إلكترونات التكافؤ} الذي
يعود إلى ذرة واحدة \emph{زوج حر}\index{زوج حر}. و\emph{بنية
لويس}\index{بنية لويس} تُظهر كل \omterm{def:b1:quantum-numbers:configuration}{إلكترون تكافؤ}
في جزيء أو أيون، روابطَ (خطوطًا بين الذرات) وأزواجًا
حرة (خطوطًا أو أزواج نقاط على الذرة). ووفق \emph{قاعدة
الثمانية}\index{قاعدة الثمانية} تحاط ذرات \omterm{def:b1:periodicity:table}{الدورة} 2 (C، N، O، F)
في البنية المستقرة بأربعة أزواج، أي ثمانية إلكترونات؛ ووفق
\emph{قاعدة الثنائية}\index{قاعدة الثنائية} يحاط الهيدروجين بزوج واحد.
\end{definition}

\begin{definition}[الشحنة الشكلية]\label{def:b1:lewis-resonance-vsepr:formal-charge}
\emph{الشحنة الشكلية}\index{شحنة شكلية} لذرة في \omterm{def:b1:lewis-resonance-vsepr:lewis}{بنية لويس}
هي
\[
  q_F = N_v - N_{\text{حر}} - \tfrac12 N_{\text{رابط}},
\]
حيث $N_v$ عدد \omterm{def:b1:quantum-numbers:configuration}{إلكترونات التكافؤ} في الذرة الحرة،
و$N_{\text{حر}}$ عدد الإلكترونات في أزواجها الحرة،
و$N_{\text{رابط}}$ عدد الإلكترونات في الروابط التي تكوّنها. وتُكتب
$\oplus$ أو $\ominus$ بجوار الذرة حين لا تكون معدومة.
\end{definition}

\begin{proposition}[مجموع الشحنات الشكلية يساوي الشحنة]\label{prop:b1:lewis-resonance-vsepr:formal-sum}
مجموع \omterm{def:b1:lewis-resonance-vsepr:formal-charge}{الشحنات الشكلية} لذرات \omterm{def:b1:lewis-resonance-vsepr:lewis}{بنية لويس} يساوي
الشحنة الكلية للجزيء أو للأيون.
\end{proposition}

\begin{proof}
بجمع $q_F$ على الذرات نحصل على $\sum N_v - (\text{إلكترونات الأزواج
الحرة} + \text{إلكترونات الروابط})$، لأن إلكترونات كل رابطة
تُعدّ نصفين على ذرتيها. وما بين القوسين هو العدد الكلي
\omterm{def:b1:quantum-numbers:configuration}{لإلكترونات التكافؤ} المرسومة في البنية، وهو يساوي $\sum N_v$ مطروحًا منه
شحنة النوع $z$. ومنه $\sum q_F = z$.
\end{proof}

\begin{method}[رسم بنية لويس]\label{met:b1:lewis-resonance-vsepr:lewis}
لرسم \omterm{def:b1:lewis-resonance-vsepr:lewis}{بنية لويس} لجزيء أو لأيون:
\begin{enumerate}
  \item عُدّ \omterm{def:b1:quantum-numbers:configuration}{إلكترونات التكافؤ} $N = \sum N_v - z$ والأزواج $N/2$؛
  \item ارسم الهيكل: الذرة الأقل كهرسلبية (غير H) هي
        المركزية عادة؛ والهيدروجين والفلور طرفيان دائمًا؛
  \item أكمل ثمانيات الذرات الطرفية \omterm{def:b1:lewis-resonance-vsepr:lewis}{بأزواج حرة}، ثم
        ضع الأزواج الباقية على الذرة المركزية؛
  \item إن نقص الذرةَ المركزيةَ ثُمانٌ، فحوّل أزواجًا حرة من
        جاراتها إلى روابط متعددة؛
  \item احسب \omterm{def:b1:lewis-resonance-vsepr:formal-charge}{الشحنات الشكلية}؛ وفضّل من بين البنى الممكنة
        ذات أقل عدد من \omterm{def:b1:lewis-resonance-vsepr:formal-charge}{الشحنات الشكلية}، والتي تقع فيها الشحنات السالبة على
        الذرات الأكثر كهرسلبية.
\end{enumerate}
\end{method}

\begin{example}[حمض النتريك]\label{ex:b1:lewis-resonance-vsepr:hno3}
\ce{HNO3}: $N = 1 + 5 + 3 \times 6 = 24$ إلكترونًا، أي 12 زوجًا. الهيكل:
النيتروجين مركزي، مرتبط بثلاث ذرات أكسجين تحمل إحداها
الهيدروجين. وإكمال الثمانيات بروابط أحادية يترك للنيتروجين ستة
إلكترونات؛ فيصير \omterm{def:b1:lewis-resonance-vsepr:lewis}{زوج حر} لإحدى ذرات الأكسجين رابطة \ce{N=O}. \omterm{def:b1:lewis-resonance-vsepr:formal-charge}{الشحنات الشكلية}:
النيتروجين $5 - 0 - 4 = +1$؛ والأكسجين المرتبط برابطة أحادية دون هيدروجين
$6 - 6 - 1 = -1$؛ والبقية 0. فالمجموع 0، كما
تقتضي \cref{prop:b1:lewis-resonance-vsepr:formal-sum}.
\end{example}

\begin{omfigure}
\setchemfig{atom sep=2.4em}
\begin{tikzpicture}[font=\small]
  \node at (0,0) {\large\chemfig{H-\charge{90=\|,270=\|}{O}-\charge{90:2pt=$\scriptstyle\oplus$}{N}(=[:45]\charge{0=\|,90=\|}{O})-[:-45]\charge{0=\|,240=\|,300=\|,45:3pt=$\scriptstyle\ominus$}{O}}};
  \node[below] at (0,-1.5) {\foreignlanguage{arabic}{حمض النتريك، \ce{HNO3}}};
  \node at (5.2,0) {\large\chemfig{\charge{180=\|}{N}~\charge{0=\|}{N}}};
  \node[below] at (5.2,-1.5) {\foreignlanguage{arabic}{ثنائي النيتروجين، \ce{N2}}};
  \node at (9.6,0) {\large\chemfig{\charge{135=\|,225=\|,90:2pt=$\scriptstyle\ominus$}{N}=\charge{90:2pt=$\scriptstyle\oplus$}{N}=\charge{45=\|,315=\|}{O}}};
  \node[below] at (9.6,-1.5) {\foreignlanguage{arabic}{أحادي أكسيد ثنائي النيتروجين، \ce{N2O}}};
\end{tikzpicture}

{\small \omterm{def:b1:lewis-resonance-vsepr:lewis}{بنى لويس} بأزواجها الحرة وشحناتها الشكلية.}
\end{omfigure}

\section{ما وراء الثمانية}

\begin{definition}[الجزيئات الناقصة الإلكترونات والفائقة التكافؤ]\label{def:b1:lewis-resonance-vsepr:hypervalent}
الجزيء الذي تحيط فيه بذرة أقل من ثمانية \omterm{def:b1:quantum-numbers:configuration}{إلكترونات تكافؤ}
\emph{ناقص الإلكترونات}\index{ناقص الإلكترونات}: ففي \ce{BF3} للبور
ستة. والجزيء الذي تحاط فيه ذرة من \omterm{def:b1:periodicity:table}{الدورة} 3 أو ما بعدها
بأكثر من أربعة أزواج \emph{فائق التكافؤ}\index{فائق التكافؤ}: ففي
\ce{PCl5} يكوّن الفوسفور خمس روابط، وفي \ce{SF6} يكوّن الكبريت ستًّا. أما الجزيئات
ذات العدد الفردي من الإلكترونات، مثل \ce{NO} و\ce{NO2}، فلا يمكنها
أن تحقق \omterm{def:b1:lewis-resonance-vsepr:lewis}{قاعدة الثمانية} على كل ذرة: فأحد الإلكترونات منفرد.
\end{definition}

\begin{remark}[لماذا الدورة 3 لا الدورة 2]\label{rem:b1:lewis-resonance-vsepr:hypervalent}
يكوّن النيتروجين \ce{NF3} ولا يكوّن \ce{NF5} أبدًا، في حين يكوّن الفوسفور
\ce{PF3} و\ce{PF5} كليهما. فذرة \omterm{def:b1:periodicity:table}{الدورة} 2 أصغر من أن ترتبط بأكثر من
أربعة جيران؛ وذرة \omterm{def:b1:periodicity:table}{الدورة} 3 كبيرة بما يكفي لتتسع لخمسة أو ستة.
ونجد رسوم \ce{SO4^{2-}} أو \ce{H2SO4} بروابط ثنائية للكبريت (12
إلكترونًا حول S) أو بروابط أحادية \omterm{def:b1:lewis-resonance-vsepr:formal-charge}{وشحنات شكلية} (ثُمانٌ)،
كلتيهما؛ والثانية أصدق تمثيلًا لتوزع الشحنة، وكلتاهما
تعطي الهندسة الصحيحة.
\end{remark}

\begin{definition}[حمض لويس، قاعدة لويس، الرابطة التناسقية]\label{def:b1:lewis-resonance-vsepr:lewis-acid}
\emph{حمض لويس}\index{حمض لويس} نوع كيميائي قادر على تقبّل
زوج إلكترونات في موضع شاغر من طبقة تكافئه (\ce{BF3}،
\ce{AlCl3}، \ce{H+}، كاتيونات الفلزات)؛ و\emph{قاعدة لويس}\index{قاعدة لويس}
نوع يملك \omterm{def:b1:lewis-resonance-vsepr:lewis}{زوجًا حرًا} يمكنه أن يشارك به (\ce{NH3}، \ce{H2O}،
\ce{F-}). وحين تمنح القاعدةُ الزوجَ للحمض، تكون الرابطة المتكوّنة
\emph{رابطة تناسقية}\index{رابطة تناسقية}: \omterm{def:b1:lewis-resonance-vsepr:lewis}{رابطة تساهمية} أتى
إلكتروناها كلاهما من أحد الشريكين.
\end{definition}

\begin{example}[الأمونيا وثلاثي فلوريد البور]\label{ex:b1:lewis-resonance-vsepr:adduct}
\ce{NH3 + BF3 -> H3NBF3}. في ناتج الإضافة يكمل البور ثُمانَه؛
ويحمل النيتروجين، الذي صار يشارك بزوجه الحر، \omterm{def:b1:lewis-resonance-vsepr:formal-charge}{الشحنة الشكلية}
$+1$ ويحمل البور $-1$. ومتى تكوّنت الرابطة \ce{N-B} صارت \omterm{def:b1:lewis-resonance-vsepr:lewis}{رابطة
تساهمية} عادية.
\end{example}

\begin{omfigure}
\setchemfig{atom sep=2.1em}
\begin{tikzpicture}[font=\small]
  \node at (0,0) {\large\chemfig{H-\charge{90=\|}{N}(-[:240]H)(-[:300]H)}};
  \node at (1.6,0) {$+$};
  \node at (3.1,0) {\large\chemfig{B(-[:90]F)(-[:210]F)(-[:330]F)}};
  \draw[-{Stealth}, thick] (4.3,0) -- (5.3,0);
  \node at (7.4,0) {\large\chemfig{H-\charge{45:1pt=$\scriptstyle\oplus$}{N}(-[:90]H)(-[:270]H)-\charge{45:1pt=$\scriptstyle\ominus$}{B}(-[:90]F)(-[:270]F)-F}};
\end{tikzpicture}

{\small \omterm{def:b1:lewis-resonance-vsepr:lewis-acid}{رابطة تناسقية}: \omterm{def:b1:lewis-resonance-vsepr:lewis}{الزوج الحر} للأمونيا (\omterm{def:b1:lewis-resonance-vsepr:lewis-acid}{قاعدة لويس}) يملأ
الموضع الشاغر لثلاثي فلوريد البور (\omterm{def:b1:lewis-resonance-vsepr:lewis-acid}{حمض لويس}). ولم تُرسم أزواج الفلور
الحرة.}
\end{omfigure}

\section{الرنين}

\begin{definition}[بنى الرنين، الهجين الرنيني]\label{def:b1:lewis-resonance-vsepr:resonance}
حين يمكن وضع إلكترونات جزيء في عدة \omterm{def:b1:lewis-resonance-vsepr:lewis}{بنى لويس}
لا تختلف إلا بمواضع روابط $\pi$ \omterm{def:b1:lewis-resonance-vsepr:lewis}{والأزواج
الحرة} --- مع بقاء الأنوية في أماكنها --- تكون كل منها \emph{بنية
رنين}\index{بنية رنين} (أو صيغة ميزوميرية)، تربطها بالأخريات
سهم مزدوج الرأس $\leftrightarrow$. والجزيء الحقيقي
ليس أيًّا منها: إنه \emph{الهجين الرنيني}\index{هجين رنيني}،
بنية واحدة توزعها الإلكتروني متوسط مرجَّح
لتوزعاتها.
\end{definition}

\begin{remark}[سهم، لا توازن]\label{rem:b1:lewis-resonance-vsepr:arrow}
لا يصف السهم $\leftrightarrow$ تفاعلًا: فالأوزون لا ينتقل
من بنية إلى أخرى. الهجين جزيء واحد،
يوصف بعدة رسوم لأن رسمًا واحدًا لا يكفي. ولا
تستعمل أبدًا $\rightleftharpoons$ بين \omterm{def:b1:lewis-resonance-vsepr:resonance}{بنى الرنين}.
\end{remark}

\begin{method}[كتابة بنى الرنين وترجيحها]\label{met:b1:lewis-resonance-vsepr:resonance}
لإيجاد \omterm{def:b1:lewis-resonance-vsepr:resonance}{بنى الرنين} لنوع كيميائي:
\begin{enumerate}
  \item انطلق من \omterm{def:b1:lewis-resonance-vsepr:lewis}{بنية لويس} واحدة؛ وابحث عن \omterm{def:b1:lewis-resonance-vsepr:lewis}{زوج حر} أو
        رابطة $\pi$ مجاورة لرابطة $\pi$ أو لموضع شاغر أو لشحنة
        موجبة؛
  \item انقل الأزواج بأسهم منحنية (يصير \omterm{def:b1:lewis-resonance-vsepr:lewis}{الزوج الحر} رابطة $\pi$، وتصير
        الرابطة $\pi$ \omterm{def:b1:lewis-resonance-vsepr:lewis}{زوجًا حرًا})، دون أن تنقل نواة أبدًا ودون أن
        تتجاوز ثُمانَ ذرة من \omterm{def:b1:periodicity:table}{الدورة} 2؛
  \item أعد حساب \omterm{def:b1:lewis-resonance-vsepr:formal-charge}{الشحنات الشكلية}؛
  \item رجّح البنى: أهمها ما يحترم \omterm{def:b1:lewis-resonance-vsepr:lewis}{قاعدة
        الثمانية}، وما فيه أقل عدد من \omterm{def:b1:lewis-resonance-vsepr:formal-charge}{الشحنات الشكلية}، وما يضع الشحنات السالبة على
        الذرات الأكثر كهرسلبية؛ والبنى المتكافئة متساوية
        الوزن.
\end{enumerate}
\end{method}

\begin{example}[الأوزون، النترات، البنزين]\label{ex:b1:lewis-resonance-vsepr:examples}
للأوزون بنيتان متكافئتان: كل رابطة \ce{O-O} ثنائية في إحداهما
وأحادية في الأخرى، فلكلتيهما رتبة رابطة $1.5$ والطول
نفسه، \qty{127.8}{pm}، بين \ce{O=O} و\ce{O-O}. ولأيون النترات
\ce{NO3-} ثلاث بنى متكافئة؛ فرتبة كل رابطة \ce{N-O}
$4/3$. وللبنزين، \ce{C6H6}، بنيتا كيكوليه؛ وروابطه \ce{C-C} الست
متساوية، \qty{139.7}{pm}، بين الرابطة الأحادية في الإيثان
(\qty{153.6}{pm}) والرابطة الثنائية في الإيثين (\qty{133.9}{pm}).
% ledger: geo:O3.rOO, geo:C6H6.rCC, geo:C2H6.rCC, geo:C2H4.rCC
\end{example}

\begin{omfigure}
\vspace*{10pt}
\large
\setchemfig{atom sep=2.4em}
\schemestart
\chemfig{\charge{135=\|,225=\|}{O}=[:-30]\charge{270=\|,90:2pt=$\scriptstyle\oplus$}{O}-[:30]\charge{90=\|,0=\|,270=\|,45:5pt=$\scriptstyle\ominus$}{O}}
\arrow{<->}
\chemfig{\charge{90=\|,180=\|,270=\|,135:5pt=$\scriptstyle\ominus$}{O}-[:-30]\charge{270=\|,90:2pt=$\scriptstyle\oplus$}{O}=[:30]\charge{45=\|,315=\|}{O}}
\schemestop

\medskip
\schemestart
\chemfig{*6(=-=-=-)}
\arrow{<->}
\chemfig{*6(-=-=-=)}
\schemestop
\hspace{3em}
\chemfig{\charge{270:2pt=$\scriptstyle\oplus$}{N}(=[:90]\charge{45=\|,135=\|}{O})(-[:210]\charge{150=\|,240=\|,300=\|,90:2pt=$\scriptstyle\ominus$}{O})-[:330]\charge{30=\|,300=\|,240=\|,90:2pt=$\scriptstyle\ominus$}{O}}
\normalsize
\par\vspace{14pt}
{\small الرنين في الأوزون (يصير \omterm{def:b1:lewis-resonance-vsepr:lewis}{زوج حر} لذرة الأكسجين اليمنى
رابطة $\pi$ بينما تصير الرابطة $\pi$ اليسرى \omterm{def:b1:lewis-resonance-vsepr:lewis}{زوجًا
حرًا}) وفي البنزين، وإحدى
بنى أيون النترات الثلاث المتكافئة، الذي تتوزع شحنته
على ذرات الأكسجين الثلاث.}
\end{omfigure}

\begin{definition}[الرابطة $\sigma$، الرابطة $\pi$]\label{def:b1:lewis-resonance-vsepr:sigma-pi}
الرابطة الأحادية \emph{رابطة $\sigma$}\index{رابطة سيغما@رابطة $\sigma$}:
زوجها الإلكتروني واقع على المحور الواصل بين النواتين، ناتج عن
التداخل المحوري لمدارين. والرابطتان الثانية والثالثة في الرابطة
المتعددة \emph{روابط $\pi$}\index{رابطة باي@رابطة $\pi$}: زوجاهما على
جانبي المحور، ناتجان عن التداخل الجانبي لمداري p
عموديين عليه. فالرابطة الثنائية رابطة $\sigma$ ورابطة $\pi$؛
والرابطة الثلاثية رابطة $\sigma$ ورابطتا $\pi$.
\end{definition}

\begin{omfigure}
\begin{tikzpicture}[font=\small]
  % sigma: two lobes overlapping head-on
  \fill[omDef!35, opacity=0.8] (0,0) ellipse (0.9 and 0.42);
  \fill[omProp!35, opacity=0.8] (1.1,0) ellipse (0.9 and 0.42);
  \fill (0.1,0) circle (1.6pt); \fill (1.0,0) circle (1.6pt);
  \draw[dashed, black!50] (-1.2,0) -- (2.3,0);
  \node[below] at (0.55,-1.05) {\foreignlanguage{arabic}{$\sigma$: تداخل محوري، على المحور}};
  % pi: two p orbitals side by side
  \begin{scope}[xshift=5.2cm]
    \foreach \x/\c in {0/omDef, 1.2/omProp} {
      \fill[\c!35, opacity=0.85] (\x,0.45) ellipse (0.32 and 0.45);
      \fill[\c!15, opacity=0.85] (\x,-0.45) ellipse (0.32 and 0.45);
      \draw[\c!70!black] (\x,0.45) ellipse (0.32 and 0.45) (\x,-0.45) ellipse (0.32 and 0.45);
      \fill (\x,0) circle (1.6pt);
    }
    \draw[dashed, black!50] (-0.7,0) -- (1.9,0);
    \draw[omThm, thick, densely dotted] (0.25,0.75) -- (0.95,0.75) (0.25,-0.75) -- (0.95,-0.75);
    \node[below] at (0.6,-1.05) {\foreignlanguage{arabic}{$\pi$: تداخل جانبي، فوق المحور وتحته}};
  \end{scope}
\end{tikzpicture}

{\small نوعا التداخل. تسمح الرابطة $\sigma$ بالدوران حول
محورها؛ ولا تسمح به الرابطة $\pi$، ولهذا كانت الرابطة الثنائية صلبة
(\cref{ch:b1:stereochemistry-in-depth}).}
\end{omfigure}

\section{نموذج VSEPR}

\begin{definition}[نموذج VSEPR، العدد الفراغي]\label{def:b1:lewis-resonance-vsepr:vsepr}
في \emph{نموذج VSEPR}\index{نموذج VSEPR} (تنافر أزواج
إلكترونات طبقة التكافؤ)، يتنافر الأزواج حول ذرة مركزية A
ويتخذ الترتيب الذي يُبقيها أبعد ما يمكن بعضها عن بعض. ويُكتب
الجزيء \ce{AX_nE_m}، حيث $n$ عدد الذرات المرتبطة
بالذرة A (تُعدّ الرابطة المتعددة مرة واحدة) و$m$ عدد \omterm{def:b1:lewis-resonance-vsepr:lewis}{الأزواج الحرة} للذرة
A. ويحدد \emph{العدد الفراغي}\index{عدد فراغي} $n + m$
ترتيب الأزواج: 2 خطي، 3 مثلثي مستوٍ، 4 رباعي الأوجه، 5
هرمي ثلاثي مزدوج، 6 ثماني الأوجه. وشكل الجزيء هو
ترتيب ذراته وحدها.
\end{definition}

\begin{proposition}[هندسات VSEPR]\label{prop:b1:lewis-resonance-vsepr:geometries}
يتنبأ \omterm{def:b1:lewis-resonance-vsepr:vsepr}{نموذج VSEPR} بالأشكال والزوايا المثالية الواردة في الجدول أدناه؛
والزوايا المقيسة قريبة منها.
\begin{center}
\small
\begin{tabular}{llllc}
\hline
النوع & ترتيب الأزواج & الشكل & مثال & الزاوية المقيسة \\
\hline
\ce{AX2} & خطي & خطي & \ce{CO2} & $180^\circ$ \\
\ce{AX3} & مثلثي مستوٍ & مثلثي مستوٍ & \ce{BF3} & $120^\circ$ \\
\ce{AX2E} & مثلثي مستوٍ & منحنٍ & \ce{SO2} & $119.5^\circ$ \\
\ce{AX4} & رباعي الأوجه & رباعي الأوجه & \ce{CH4} & $109.5^\circ$ \\
\ce{AX3E} & رباعي الأوجه & هرمي ثلاثي & \ce{NH3} & $106.7^\circ$ \\
\ce{AX2E2} & رباعي الأوجه & منحنٍ & \ce{H2O} & $104.5^\circ$ \\
\ce{AX5} & هرمي ثلاثي مزدوج & هرمي ثلاثي مزدوج & \ce{PCl5} & $90^\circ$، $120^\circ$ \\
\ce{AX4E} & هرمي ثلاثي مزدوج & أرجوحي & \ce{SF4} & $101.6^\circ$، $173.1^\circ$ \\
\ce{AX3E2} & هرمي ثلاثي مزدوج & على شكل حرف T & \ce{ClF3} & $87.5^\circ$ \\
\ce{AX2E3} & هرمي ثلاثي مزدوج & خطي & \ce{XeF2} & $180^\circ$ \\
\ce{AX6} & ثماني الأوجه & ثماني الأوجه & \ce{SF6} & $90^\circ$ \\
\ce{AX5E} & ثماني الأوجه & هرمي مربع القاعدة & \ce{BrF5} & --- \\
\ce{AX4E2} & ثماني الأوجه & مربع مستوٍ & \ce{XeF4} & --- \\
\hline
\end{tabular}
\end{center}
% ledger: geo:CO2.aOCO, geo:BF3.aFBF, geo:SO2.aOSO, geo:CH4.aHCH,
% geo:NH3.aHNH, geo:H2O.aHOH, geo:SF6.aFSF, geo:SF4.aFSF2, geo:SF4.aFSF,
% geo:ClF3.aFClF, geo:XeF2.aFXeF
\end{proposition}
\admitted

\begin{proposition}[الأزواج الحرة تُضيّق الزوايا]\label{prop:b1:lewis-resonance-vsepr:lone-pair-angles}
\omterm{def:b1:lewis-resonance-vsepr:lewis}{الزوج الحر}، الذي تمسكه نواة واحدة، ينتشر أكثر من \omterm{def:b1:lewis-resonance-vsepr:lewis}{الزوج الرابط}
ويتنافر بقوة أكبر: تتناقص التنافرات بالترتيب
\babelsublr{\foreignlanguage{arabic}{\omterm{def:b1:lewis-resonance-vsepr:lewis}{زوج حر}/\omterm{def:b1:lewis-resonance-vsepr:lewis}{زوج حر}} $>$ \foreignlanguage{arabic}{\omterm{def:b1:lewis-resonance-vsepr:lewis}{زوج حر}/\omterm{def:b1:lewis-resonance-vsepr:lewis}{زوج رابط}} $>$ \foreignlanguage{arabic}{\omterm{def:b1:lewis-resonance-vsepr:lewis}{زوج رابط}/\omterm{def:b1:lewis-resonance-vsepr:lewis}{زوج رابط}}}.
ولذلك تضيق الزوايا بين الروابط كلما أضيفت \omterm{def:b1:lewis-resonance-vsepr:lewis}{أزواج حرة}
($109.5^\circ$ في \ce{CH4}، و$106.7^\circ$ في \ce{NH3}، و$104.5^\circ$ في
\ce{H2O})، وفي الهرم الثلاثي المزدوج تحتل \omterm{def:b1:lewis-resonance-vsepr:lewis}{الأزواج الحرة} المواضع
الاستوائية، حيث لا يكون لها إلا جاران عند $90^\circ$.
% ledger: geo:CH4.aHCH, geo:NH3.aHNH, geo:H2O.aHOH
\end{proposition}
\admitted

\begin{omfigure}
\setchemfig{atom sep=2.6em}
\renewcommand{\arraystretch}{1.2}
\begin{tabular}{ccccc}
\chemfig{O=C=O} &
\chemfig{B(-[:90]F)(-[:210]F)(-[:330]F)} &
\chemfig{C(-[:90]H)(-[:200]H)(<[:280]H)(<:[:335]H)} &
\chemfig{P(-[:90]Cl)(-[:270]Cl)(-[:180]Cl)(<[:320]Cl)(<:[:30]Cl)} &
\chemfig{S(-[:90]F)(-[:270]F)(-[:180]F)(-[:0]F)(<[:225]F)(<:[:45]F)} \\
\small \ce{AX2}, \ce{CO2} & \small \ce{AX3}, \ce{BF3} & \small \ce{AX4}, \ce{CH4} &
\small \ce{AX5}, \ce{PCl5} & \small \ce{AX6}, \ce{SF6} \\[16pt]
\chemfig{\charge{90=\:}{S}(=[:210]O)(=[:330]O)} &
\chemfig{\charge{90=\:}{N}(-[:200]H)(<[:280]H)(<:[:335]H)} &
\chemfig{\charge{60=\:,120=\:}{O}(-[:215]H)(-[:325]H)} &
\chemfig{\charge{0=\:}{S}(-[:90]F)(-[:270]F)(<[:210]F)(<:[:150]F)} &
\chemfig{Xe(-[:0]F)(-[:90]F)(-[:180]F)(-[:270]F)} \\
\small \ce{AX2E}, \ce{SO2} & \small \ce{AX3E}, \ce{NH3} & \small \ce{AX2E2}, \ce{H2O} &
\small \ce{AX4E}, \ce{SF4} & \small \ce{AX4E2}, \ce{XeF4} \\
\end{tabular}

\medskip
{\small الأشكال التي يتنبأ بها \omterm{def:b1:lewis-resonance-vsepr:vsepr}{نموذج VSEPR}. يتجه الإسفين نحو القارئ،
والإسفين المخطط بعيدًا عنه؛ والروابط العادية تقع في مستوي الصفحة. \omterm{def:b1:lewis-resonance-vsepr:lewis}{والأزواج الحرة}
للذرة المركزية مرسومة أزواج نقاط؛ و\ce{XeF4} مرسوم من
الأمام، وزوجاه الحران يتجهان نحو القارئ وبعيدًا عنه.}
\end{omfigure}

\begin{method}[التنبؤ بالشكل]\label{met:b1:lewis-resonance-vsepr:vsepr}
للتنبؤ بالشكل حول ذرة مركزية:
\begin{enumerate}
  \item ارسم \omterm{def:b1:lewis-resonance-vsepr:lewis}{بنية لويس} (\cref{met:b1:lewis-resonance-vsepr:lewis})؛
  \item عُدّ الذرات المرتبطة بالذرة المركزية ($n$) وأزواجها
        الحرة ($m$)؛ واكتب \ce{AX_nE_m}؛
  \item اقرأ الترتيب من $n + m$ والشكل من
        مواضع الذرات $n$؛
  \item صحّح الزوايا المثالية: تشغل \omterm{def:b1:lewis-resonance-vsepr:lewis}{الأزواج الحرة} والروابط المتعددة حيزًا
        أكبر مما تشغله الروابط الأحادية.
\end{enumerate}
\end{method}

\begin{definition}[التهجين]\label{def:b1:lewis-resonance-vsepr:hybridisation}
\emph{تهجين}\index{تهجين} الذرة وصفٌ
لروابطها بمدارات تتجه على امتداد الاتجاهات
التي يتنبأ بها \omterm{def:b1:lewis-resonance-vsepr:vsepr}{نموذج VSEPR}: فالذرة ذات \omterm{def:b1:lewis-resonance-vsepr:vsepr}{العدد الفراغي} 4 يقال إنها \babelsublr{sp$^3$}
(أربعة مدارات متكافئة عند $109.5^\circ$)، وذات \omterm{def:b1:lewis-resonance-vsepr:vsepr}{العدد الفراغي} 3 \babelsublr{sp$^2$}
(ثلاثة مدارات عند $120^\circ$ في مستوٍ، مع مدار p باقٍ
عمودي عليه)، وذات \omterm{def:b1:lewis-resonance-vsepr:vsepr}{العدد الفراغي} 2 sp (مداران عند $180^\circ$،
مع مداري p باقيين). ومدارات p الباقية هي التي تكوّن روابط $\pi$.
\end{definition}

\begin{example}[كربونات الكيمياء العضوية]\label{ex:b1:lewis-resonance-vsepr:carbon}
في الإيثان كل كربون \babelsublr{sp$^3$}، رباعي الأوجه؛ وفي الإيثين \babelsublr{sp$^2$}، والذرات الست
في مستوٍ واحد بزوايا قريبة من $120^\circ$ (المقيسة \ce{H-C-H}
$117.6^\circ$)؛ وفي الإيثاين sp، والذرات الأربع على مستقيم واحد. وتقصر
أطوال الروابط من 153.6 إلى 133.9 إلى \qty{120.3}{pm} كلما أضيفت روابط $\pi$.
% ledger: geo:C2H4.aHCH, geo:C2H6.rCC, geo:C2H4.rCC, geo:C2H2.rCC
\end{example}

\begin{remark}[وصف لا سبب]\label{rem:b1:lewis-resonance-vsepr:hybrid-limit}
يصف \omterm{def:b1:lewis-resonance-vsepr:hybridisation}{التهجين} هندسة معروفة سلفًا؛ ولا يفسرها،
ويخفق في طاقات الإلكترونات، التي تفسرها نظرية المدارات الجزيئية
في مجلد السنة الثانية. لكنه يبقى اللغة اليومية
للكيمياء العضوية: فعبارة «كربون \babelsublr{sp$^2$}» تعني كربونًا مستويًا
له مدار p حر لرابطة $\pi$.
\end{remark}

\section{الجزيئات القطبية: عزوم ثنائي القطب}

\begin{definition}[ثنائي قطب الرابطة، عزم ثنائي القطب]\label{def:b1:lewis-resonance-vsepr:dipole}
في رابطة بين ذرتين مختلفتي \omterm{def:b1:periodicity:electronegativity}{الكهرسلبية} تميل الإلكترونات
نحو الذرة الأكثر كهرسلبية، فتحمل شحنة جزئية
$-\delta e$ بينما يحمل شريكها $+\delta e$ ($0 < \delta < 1$).
و\emph{ثنائي قطب الرابطة}\index{ثنائي قطب الرابطة} هو المتجهة $\vec\mu = \delta e\,
\vec d$، التي طويلتها $\delta e d$، المتجهة من الشحنة الجزئية السالبة إلى الموجبة،
حيث $\vec d$ المتجهة بينهما.
و\emph{عزم ثنائي القطب}\index{عزم ثنائي القطب} لجزيء هو المجموع المتجهي
لثنائيات أقطاب روابطه (ولإسهامات أزواجه الحرة)؛ ويُعبَّر عنه
بالكولوم متر أو بالديباي، $\qty{1}{\debye} =
\qty{3.336e-30}{C.m}$. والجزيء ذو عزم ثنائي القطب غير المعدوم
\emph{جزيء قطبي}\index{جزيء قطبي}.
% ledger: const:c (1 D = 1e-21/c C.m)
\end{definition}

\begin{remark}[اصطلاح الاتجاه]\label{rem:b1:lewis-resonance-vsepr:convention}
ترسم الفيزياء متجهة ثنائي القطب من الشحنة السالبة إلى الشحنة
الموجبة، وهو الاصطلاح المتبع هنا. أما كتب الكيمياء الأقدم فترسم سهمًا
ذيله مشطوب يتجه نحو الطرف السالب؛ ولا يختلف إلا
الاتجاه.
\end{remark}

\begin{proposition}[ثنائي قطب جزيء متناظر]\label{prop:b1:lewis-resonance-vsepr:dipole-sum}
للجزيء \ce{AX_n} الذي لا \omterm{def:b1:lewis-resonance-vsepr:lewis}{أزواج حرة} على ذرته A (\ce{AX2} الخطي،
و\ce{AX3} المثلثي، و\ce{AX4} الرباعي الأوجه، \dots) \omterm{def:b1:lewis-resonance-vsepr:dipole}{عزم ثنائي قطب} معدوم، مهما
كانت قطبية روابطه. وللجزيء \ce{AX2} المنحني ذي الزاوية $\theta$
الذي لكل من رابطتيه ثنائي القطب $\mu_b$ \omterm{def:b1:lewis-resonance-vsepr:dipole}{عزم ثنائي القطب}
\[
  \mu = 2\mu_b \cos\frac{\theta}{2} .
\]
\end{proposition}

\begin{proof}
في الأشكال المتناظرة تكون متجهات الروابط متساوية الطويلة
ومرتبة اتجاهاتها ترتيبًا متناظرًا حول A، فمجموعها معدوم
(في \ce{AX2} الخطي هي متعاكسة؛ وفي \ce{AX3} عند $120^\circ$
مجموع ثلاث متجهات متساوية الطويلة معدوم؛ وفي \ce{AX4} لا يتغير المجموع
بدورانات رباعي الأوجه، فهو معدوم). وفي
الجزيء المنحني، تصنع كل متجهة رابطة الزاوية $\theta/2$ مع
المنصِّف؛ فتتلاشى المركّبات العمودية على المنصِّف وتُجمع
المركّبات على امتداده: $2\mu_b\cos(\theta/2)$.
\end{proof}

\begin{omfigure}
\begin{tikzpicture}[font=\small, >={Stealth}]
  % water
  \node[aO] (o) at (0,0) {};
  \node[aH] (h1) at (-0.95,-0.73) {};
  \node[aH] (h2) at (0.95,-0.73) {};
  \begin{scope}[on background layer]
    \draw[bond] (o.center) -- (h1.center); \draw[bond] (o.center) -- (h2.center);
  \end{scope}
  \draw[->, omProp, thick] (-0.25,0.15) -- (-0.95,-0.38);
  \draw[->, omProp, thick] (0.25,0.15) -- (0.95,-0.38);
  \draw[->, omThm, very thick] (0,-0.25) -- (0,-1.45);
  \node[omThm, right] at (0.05,-1.25) {$\vec\mu$};
  \node at (0,0.55) {$\delta^-$};
  \node at (-1.25,-1.05) {$\delta^+$}; \node at (1.25,-1.05) {$\delta^+$};
  \node[below] at (0,-1.7) {\ce{H2O}: $\mu = \qty{1.86}{\debye}$};
  % carbon dioxide
  \begin{scope}[xshift=5cm]
    \node[aO] (a) at (-1.1,0) {}; \node[aC] (c) at (0,0) {}; \node[aO] (b) at (1.1,0) {};
    \begin{scope}[on background layer]
      \foreach \d in {-1.6,1.6} {
        \draw[bond, line width=1.1pt] ([yshift=\d pt]a.center) -- ([yshift=\d pt]c.center);
        \draw[bond, line width=1.1pt] ([yshift=\d pt]c.center) -- ([yshift=\d pt]b.center);}
    \end{scope}
    \draw[->, omProp, thick] (-0.75,0.45) -- (-0.2,0.45);
    \draw[->, omProp, thick] (0.75,0.45) -- (0.2,0.45);
    \node[below] at (0,-1.7) {\foreignlanguage{arabic}{\ce{CO2}: تتلاشى ثنائيات أقطاب الروابط، $\mu = 0$}};
  \end{scope}
\end{tikzpicture}

{\small ثنائيات أقطاب الروابط (بالبرتقالي، من الشحنة الجزئية السالبة إلى
الشحنة الجزئية الموجبة) ومجموعها (بالأحمر). في الماء المنحني تتجمع؛ وفي ثنائي أكسيد
الكربون الخطي تتلاشى.}
% ledger: dip:H2O
\end{omfigure}

\begin{example}[جزيئات قطبية وغير قطبية]\label{ex:b1:lewis-resonance-vsepr:dipoles}
\omterm{def:b1:lewis-resonance-vsepr:dipole}{عزوم ثنائي القطب} للجزيئات \ce{CO2} و\ce{BF3} و\ce{CH4} و\ce{CCl4} معدومة.
أما \ce{H2O} (\qty{1.86}{\debye}) و\ce{NH3} (\qty{1.48}{\debye})
و\ce{SO2} (\qty{1.63}{\debye}) فقطبية. وعلى امتداد \ce{CH3Cl} و\ce{CH2Cl2}
و\ce{CHCl3} يهبط العزم، 1.87 ثم 1.62 ثم \qty{1.04}{\debye}، حتى ينعدم في
\ce{CCl4}: إذ تتلاشى ثنائيات أقطاب الروابط \ce{C-Cl} أكثر فأكثر. وفي
هاليدات الهيدروجين يتبع العزمُ فرقَ \omterm{def:b1:periodicity:electronegativity}{الكهرسلبية}: HF 1.83،
وHCl 1.09، وHBr 0.83، وHI \qty{0.45}{\debye}.
% ledger: dip:H2O, dip:NH3, dip:SO2, dip:CH3Cl, dip:CH2Cl2, dip:CHCl3,
% dip:CCl4, dip:HF, dip:HCl, dip:HBr, dip:HI
\end{example}

\section{تمارين}

\begin{exercise}[$\star$]\label{exo:b1:lewis-resonance-vsepr:1}
ارسم \omterm{def:b1:lewis-resonance-vsepr:lewis}{بنى لويس} للأنواع \ce{H2O} و\ce{NH3} و\ce{CO2} و\ce{HCN}
و\ce{CH2O} (الميثانال) و\ce{NH4+}، وأعطِ \omterm{def:b1:lewis-resonance-vsepr:formal-charge}{الشحنات الشكلية}.
\end{exercise}

\begin{exercise}[$\star$]\label{exo:b1:lewis-resonance-vsepr:2}
أعطِ النوع \ce{AX_nE_m} والشكل والزاوية التقريبية لكل من
\ce{H2S} و\ce{PCl3} و\ce{BF3} و\ce{SiH4} و\ce{CO3^{2-}} (الذرة المركزية
أولًا في كل صيغة).
\end{exercise}

\begin{exercise}[$\star$]\label{exo:b1:lewis-resonance-vsepr:3}
أي هذه الجزيئات قطبي: \ce{CCl4}، \ce{CH2Cl2}، \ce{BF3}،
\ce{NH3}، \ce{CO2}، \ce{SO2}؟ علّل انطلاقًا من الأشكال.
\end{exercise}

\begin{exercise}[$\star$]\label{exo:b1:lewis-resonance-vsepr:4}
عيّن \omterm{def:b1:lewis-resonance-vsepr:lewis-acid}{حمض لويس} \omterm{def:b1:lewis-resonance-vsepr:lewis-acid}{وقاعدة لويس} في \ce{H+ + H2O -> H3O+}
وفي \ce{Cu^{2+} + 4NH3 -> [Cu(NH3)4]^{2+}}. أي الروابط تناسقية؟
\end{exercise}

\begin{exercise}[$\star\star$]\label{exo:b1:lewis-resonance-vsepr:5}
اكتب \omterm{def:b1:lewis-resonance-vsepr:resonance}{بنيتي رنين} لأيون الإيثانوات \ce{CH3COO-}. بماذا
تتنبآن لطولي الرابطتين \ce{C-O}؟ قارن بحمض
الميثانويك \ce{HCOOH}، الذي تبلغ رابطتاه \ce{C-O} 120.2
و\qty{134.3}{pm}.
% ledger: geo:H2CO2.rCO, geo:H2CO2.rCO2
\end{exercise}

\begin{exercise}[$\star\star$]\label{exo:b1:lewis-resonance-vsepr:6}
ارسم \omterm{def:b1:lewis-resonance-vsepr:lewis}{بنية لويس} لأيون الكبريتات \ce{SO4^{2-}} بأربع
روابط \ce{S-O} أحادية، واحسب \omterm{def:b1:lewis-resonance-vsepr:formal-charge}{الشحنات الشكلية}، وتنبأ
بالشكل. كم \omterm{def:b1:lewis-resonance-vsepr:resonance}{بنية رنين} ذات رابطتين \ce{S=O} ودون
شحنة على الكبريت يمكن كتابتها؟
\end{exercise}

\begin{exercise}[$\star\star$]\label{exo:b1:lewis-resonance-vsepr:7}
استعمل \omterm{def:b1:lewis-resonance-vsepr:vsepr}{نموذج VSEPR} للتنبؤ بأشكال \ce{SF4} و\ce{ClF3} و\ce{XeF2}
و\ce{XeF4}. وفسّر لماذا تقع \omterm{def:b1:lewis-resonance-vsepr:lewis}{الأزواج الحرة} للهرم الثلاثي المزدوج في
المستوي الاستوائي.
\end{exercise}

\begin{exercise}[$\star\star$]\label{exo:b1:lewis-resonance-vsepr:8}
الزوايا \ce{H-X-H} هي $104.5^\circ$ في \ce{H2O} و$92.1^\circ$ في
\ce{H2S}؛ و$106.7^\circ$ في \ce{NH3} و$93.3^\circ$ في \ce{PH3}. اقترح
تفسيرًا مبنيًا على حجم الذرة المركزية
وكهرسلبيتها.
% ledger: geo:H2O.aHOH, geo:H2S.aHSH, geo:NH3.aHNH, geo:PH3.aHPH
\end{exercise}

\begin{exercise}[$\star\star$]\label{exo:b1:lewis-resonance-vsepr:9}
أعطِ \omterm{def:b1:lewis-resonance-vsepr:hybridisation}{تهجين} كل كربون \omterm{def:b1:lewis-resonance-vsepr:hybridisation}{وتهجين} النيتروجين في
\ce{CH3-C#N} (الإيثاننتريل) وفي \ce{CH2=CH-CH3}، وعدد روابط
$\sigma$ و$\pi$ في كل جزيء.
\end{exercise}

\begin{exercise}[$\star\star\star$]\label{exo:b1:lewis-resonance-vsepr:10}
\omterm{def:b1:lewis-resonance-vsepr:dipole}{عزم ثنائي القطب} للجزيء \ce{NF3} لا يتجاوز \qty{0.24}{\debye}، مقابل
\qty{1.48}{\debye} للجزيء \ce{NH3}، مع أن الرابطة \ce{N-F} أكثر قطبية
من الرابطة \ce{N-H}. فسّر الفرق مستعينًا بثنائيات أقطاب الروابط وبالزوج الحر
للنيتروجين.
% ledger: dip:NF3, dip:NH3
\end{exercise}

\begin{exercise}[$\star\star\star$]\label{exo:b1:lewis-resonance-vsepr:11}
لجزيء \ce{SO2} زاوية $119.5^\circ$ \omterm{def:b1:lewis-resonance-vsepr:dipole}{وعزم ثنائي قطب}
\qty{1.63}{\debye}. احسب \omterm{def:b1:lewis-resonance-vsepr:dipole}{ثنائي قطب رابطة} \ce{S-O} واحدة؛ ثم، مع
طول الرابطة \qty{143.2}{pm}، استنتج الشحنة الجزئية $\delta$ على كل
ذرة أكسجين.
% ledger: geo:SO2.aOSO, geo:SO2.rSO, dip:SO2
\end{exercise}

\begin{exercise}[$\star\star\star$]\label{exo:b1:lewis-resonance-vsepr:12}
ارسم أهم \omterm{def:b1:lewis-resonance-vsepr:lewis}{بنى لويس} لأحادي أكسيد ثنائي النيتروجين \ce{N2O}
(خطي، \ce{N-N-O}) واحسب \omterm{def:b1:lewis-resonance-vsepr:formal-charge}{الشحنات الشكلية}. أطوال الروابط
المقيسة \ce{N-N} \qty{112.8}{pm} و\ce{N-O} \qty{118.4}{pm}؛ وفي
\ce{N2} تبلغ الرابطة الثلاثية \qty{109.8}{pm}. أي البنيتين أثقل
وزنًا؟
% ledger: geo:N2O.rNN, geo:N2O.rNO, re:N2
\end{exercise}

\section{مسألة: جزيئات صاعقة البرق}

\begin{problem}[{مسألة نهاية الأسبوع --- الأنواع النيتروجينية والأكسجينية التي تنشأ
حين يسخّن البرق الهواء: بنى لويس، والرنين، والأشكال، والشحنات
الجزئية للماء}]
\label{pb:b1:lewis-resonance-vsepr:1}
تسخّن صاعقة البرق الهواء إلى آلاف الدرجات: فيعطي \ce{N2} و\ce{O2}
\ce{NO}، ثم \ce{NO2} و\ce{N2O} والأوزون وأخيرًا حمض النتريك في
المطر. معطيات: $\qty{1}{\debye} = \qty{3.336e-30}{C.m}$، $e =
\qty{1.602e-19}{C}$. وتُعطى القيم المقيسة حيث تلزم.

\medskip
\noindent\textbf{الجزء الأول --- \omterm{def:b1:lewis-resonance-vsepr:lewis}{بنى لويس}.}
\begin{enumerate}
  \item عُدّ \omterm{def:b1:quantum-numbers:configuration}{إلكترونات التكافؤ} في \ce{N2} و\ce{NO} و\ce{NO2}
        و\ce{N2O} و\ce{O3} و\ce{HNO3}.
  \item ارسم \omterm{def:b1:lewis-resonance-vsepr:lewis}{بنية لويس} للجزيء \ce{N2}.
  \item ارسم \omterm{def:b1:lewis-resonance-vsepr:lewis}{بنية لويس} للجزيء \ce{NO}. لماذا لا يمكن تحقيق \omterm{def:b1:lewis-resonance-vsepr:lewis}{قاعدة الثمانية}
        على الذرتين كلتيهما؟
  \item ارسم \omterm{def:b1:lewis-resonance-vsepr:lewis}{بنيتي لويس} للجزيء \ce{N2O}، الأولى \ce{N=N=O} والثانية
        \ce{N#N-O}، وأعطِ \omterm{def:b1:lewis-resonance-vsepr:formal-charge}{الشحنات الشكلية}.
  \item ارسم \omterm{def:b1:lewis-resonance-vsepr:lewis}{بنية لويس} للأوزون وأعطِ \omterm{def:b1:lewis-resonance-vsepr:formal-charge}{الشحنات الشكلية}.
  \item ارسم \omterm{def:b1:lewis-resonance-vsepr:lewis}{بنية لويس} للجزيء \ce{HNO3} وأعطِ \omterm{def:b1:lewis-resonance-vsepr:formal-charge}{الشحنات الشكلية}.
  \item أي الأنواع الستة لها \omterm{def:b1:quantum-numbers:unpaired}{إلكترون منفرد}؟
\end{enumerate}

\medskip
\noindent\textbf{الجزء الثاني --- الرنين وأطوال الروابط.}
\begin{enumerate}[resume]
  \item اكتب \omterm{def:b1:lewis-resonance-vsepr:resonance}{بنيتي الرنين} للأوزون. بأي رتبة رابطة
        تتنبآن؟
  \item تبلغ الرابطة \ce{O-O} \qty{127.8}{pm} في الأوزون،
        و\qty{120.8}{pm} في \ce{O2} و\qty{147.5}{pm} في \ce{H2O2}. هل
        يتأكد التنبؤ؟
  \item اكتب \omterm{def:b1:lewis-resonance-vsepr:resonance}{بنى الرنين} الثلاث للأيون \ce{NO3-} وأعطِ
        رتبة كل رابطة \ce{N-O}.
  \item في \ce{HNO3} تبلغ الروابط \ce{N-O} الثلاث 140.6 و121.1
        و\qty{119.9}{pm}. انسب كلًّا منها إلى رابطتها وفسّر مستعينًا بالرنين.
  \item فسّر لماذا كانت الزاوية \ce{O-N-O} في \ce{NO2} ($134^\circ$)
        أكبر من $120^\circ$.
  \item أي البنيتين في السؤال 4 تضع \omterm{def:b1:lewis-resonance-vsepr:formal-charge}{الشحنة الشكلية}
        السالبة على الذرة الأكثر كهرسلبية؟
\end{enumerate}

\medskip
\noindent\textbf{الجزء الثالث --- الأشكال والقطبية.}
\begin{enumerate}[resume]
  \item أعطِ النوع \ce{AX_nE_m} للذرة المركزية في \ce{O3}
        و\ce{N2O} و\ce{NO3-} ولنيتروجين \ce{HNO3}، وأشكالها.
  \item زاوية الأوزون $116.8^\circ$. فسّر لماذا هي أدنى من
        $120^\circ$.
  \item للأوزون \omterm{def:b1:lewis-resonance-vsepr:dipole}{عزم ثنائي قطب} \qty{0.53}{\debye}، وللجزيء \ce{N2O}
        \qty{0.16}{\debye}، وليس للجزيء \ce{N2} أي عزم. فسّر.
  \item لماذا كان \ce{CO2} غير قطبي بينما \ce{SO2} قطبي؟
  \item تنبأ بزاوية \ce{NH3} مقارنةً بالقيمة $109.5^\circ$، ثم
        قارن بالقيمة المقيسة $106.7^\circ$.
  \item السؤال نفسه للماء ($104.5^\circ$).
\end{enumerate}
% ledger: geo:O3.rOO, re:O2, geo:H2O2.rOO, geo:HNO3.rNO, geo:HNO3.rNO2,
% geo:HNO3.rNO3, geo:NO2.aONO, geo:O3.aOOO, dip:O3, dip:N2O

\medskip
\noindent\textbf{الجزء الرابع --- الشحنات الجزئية للماء.}
\omterm{def:b1:lewis-resonance-vsepr:dipole}{عزم ثنائي القطب} للماء \qty{1.857}{\debye}، وزاويته
$104.48^\circ$، وطول رابطته \ce{O-H} \qty{95.8}{pm}.
% ledger: dip:H2O, geo:H2O.aHOH, geo:H2O.rOH
\begin{enumerate}[resume]
  \item عبّر عن \omterm{def:b1:lewis-resonance-vsepr:dipole}{عزم ثنائي القطب} للماء بدلالة \omterm{def:b1:lewis-resonance-vsepr:dipole}{ثنائي قطب
        الرابطة} $\mu_{OH}$ والزاوية.
  \item احسب $\mu_{OH}$ بالديباي.
  \item حوّله إلى الكولوم متر.
  \item لو كان \omterm{def:b1:lewis-resonance-vsepr:dipole}{ثنائي قطب الرابطة} مكوّنًا من شحنتين تامتين $\pm e$ عند
        النواتين، فكم تكون قيمته؟
  \item استنتج الطابع الأيوني الجزئي $\delta$ للرابطة \ce{O-H}،
        وهو نسبة \omterm{def:b1:lewis-resonance-vsepr:dipole}{ثنائي قطب الرابطة} المقيس إلى ثنائي القطب ذي
        الشحنات التامة.
\end{enumerate}
\end{problem}
